\documentclass[11pt]{article}

\usepackage[margin=1.05in]{geometry}
\usepackage{amsmath,amssymb,amsthm,mathtools}
\usepackage[T1]{fontenc}
\usepackage{lmodern}
\usepackage{microtype}
\usepackage{xcolor}
\usepackage{hyperref}
\hypersetup{
  colorlinks=true,
  linkcolor=blue!55!black,
  citecolor=blue!55!black,
  urlcolor=blue!55!black,
  pdftitle={Full symmetric monodromy throughout Gallagher's weighted-lift family}
}

\newtheorem{theorem}{Theorem}
\newtheorem{lemma}[theorem]{Lemma}
\newtheorem{corollary}[theorem]{Corollary}
\theoremstyle{remark}
\newtheorem{remark}[theorem]{Remark}
\newcommand{\A}{\mathbb A}
\newcommand{\C}{\mathbb C}
\newcommand{\Mon}{\operatorname{Mon}}
\newcommand{\Gal}{\operatorname{Gal}}

\title{\bfseries Full symmetric monodromy throughout\\
Gallagher's weighted-lift family\\[3pt]
\large An attributed Brink--Gallagher corollary}
\author{Alec Kriebel, with substantial AI assistance}
\date{Research record: 25 July 2026\\
First public release: 25 July 2026, 20:15:16 UTC}

\begin{document}
\maketitle

\begin{center}
\fcolorbox{black!45}{yellow!8}{%
\begin{minipage}{0.90\textwidth}
\small
\textbf{Status, priority, and review notice.}
Gallagher has priority for the weighted-lift maps and their generic degrees.
Brink's 2004 theorem supplies the abstract two-free-coefficient Galois
calculation.  A MathOverflow answer and linked Note~19 had already proved
full symmetric monodromy for Gallagher's canonical tower in generic degrees
$3$ through $13$.  The contribution recorded here is only the uniform
all-degree, all-admissible-seed corollary obtained by combining Gallagher's
exact inverse equation with Brink's theorem.  A source audit found no earlier
statement of that combination, but cannot guarantee worldwide priority.
This work is not peer reviewed.  Exact computations check encoded algebra;
they are evidence against transcription errors, not peer review.
\end{minipage}}
\end{center}

\begin{abstract}
Gallagher constructed dimension-three Keller counterexamples of every
generic degree $d\geq3$ by a weighted-lift procedure depending on an
admissible polynomial seed.  We observe that the inverse equation for every
such seed is
\[
 \Phi(T)-PT+cQ=0,
\]
where $\Phi$ has degree $d$ and $P,Q$ are independent target parameters.
Brink's classical theorem on a polynomial with free linear and constant
coefficients therefore applies verbatim.  Exact rational recovery identifies
the root field with the function-field extension of the Keller map and rules
out extraneous denominator branches.  Consequently every admissible
weighted lift of generic degree $d$ has geometric monodromy $S_d$.
In particular every symmetric group $S_d$, $d\geq3$, is realized as the
geometric monodromy group of a Keller counterexample
$\A^3_\C\to\A^3_\C$.
\end{abstract}

\section{Statement}

For a dominant polynomial map
$F:\A^3_\C\to\A^3_\C$, write
\[
 K=\C(F_1,F_2,F_3),\qquad L=\C(x,y,z).
\]
Its generic degree is $[L:K]$, and its geometric monodromy group is the
Galois group of the Galois closure of $L/K$, in its natural action on the
generic fibre.

\begin{theorem}\label{thm:main}
Let $p\in\C[w]$ be any admissible seed in Gallagher's weighted-lift
construction and suppose $\deg p=d-1\geq2$.  The associated Keller
counterexample has generic degree $d$ and geometric monodromy group $S_d$.
\end{theorem}

\begin{corollary}\label{cor:realization}
For every $d\geq3$, the symmetric group $S_d$ occurs as the geometric
monodromy group of a Keller counterexample
$\A^3_\C\to\A^3_\C$ of generic degree $d$.
\end{corollary}

\section{The weighted lift}

Fix $b,c\in\C^\times$.  An admissible seed is a polynomial $p$ satisfying
\begin{equation}
 p(0)=0,\qquad p(1)=-c,\qquad \int_0^1p(s)\,ds=0,
 \qquad \kappa:=\frac{p'(1)}c\ne-2.                 \label{eq:seed}
\end{equation}
Define
\begin{equation}
 q(0)=0,\qquad q'(w)=\frac{w p'(w)}c,\qquad
 a=-\frac{1+\kappa}{2+\kappa}.                     \label{eq:q}
\end{equation}
Put
\[
 v=xy,\quad t=x^2z,\quad u=1+v,\quad
 \gamma=1+av+bt,\quad w=u\gamma
\]
and
\[
 \beta=c+\frac{p(w)}{\gamma},\qquad
 \alpha=u+\frac{q(w)}{\gamma^2}.
\]
Gallagher's weighted lift is
\begin{equation}
 F_p=(A,B,C)=
 \left(\frac{\alpha}{x^2},\frac{\beta}{x},x\gamma\right).
                                                               \label{eq:map}
\end{equation}

For completeness, the apparent divisions in \eqref{eq:map} cancel.
The condition $p(0)=0$ makes $p(u\gamma)/\gamma$ polynomial in $u,\gamma$.
Since $q(0)=q'(0)=0$, the same holds for
$q(u\gamma)/\gamma^2$.  At $(v,t)=(0,0)$ one has $u=\gamma=w=1$.
Integration by parts gives
\begin{equation}
 q(w)=\frac{wp(w)-\Phi(w)}c,\qquad
 \Phi(w):=\int_0^w p(s)\,ds,                       \label{eq:phi}
\end{equation}
so \eqref{eq:seed} gives $q(1)=-1$.  Thus
$\alpha(0,0)=\beta(0,0)=0$, while
\[
 \frac{\partial\alpha}{\partial v}(0,0)
 =1+\kappa(1+a)+2a=0.
\]
After substituting $v=xy$ and $t=x^2z$, $\beta$ is divisible by $x$ and
$\alpha$ by $x^2$.

We also recall the determinant calculation.  Set
\begin{equation}
 P=BC=c\gamma+p(w),\qquad
 Q=AC^2=w\gamma+q(w).                              \label{eq:PQ}
\end{equation}
Then
\[
 \det\frac{\partial(P,Q)}{\partial(w,\gamma)}=-c\gamma,
 \qquad
 \det\frac{\partial(w,\gamma)}{\partial(v,t)}=b\gamma.
\]
Moreover,
\[
 \det\frac{\partial(x,v,t)}{\partial(x,y,z)}=x^3,
 \qquad
 \det\frac{\partial(A,B,C)}{\partial(P,Q,C)}=-C^{-3}.
\]
Multiplication, using $C=x\gamma$, gives
\begin{equation}
 \det JF_p=bc\ne0.                                  \label{eq:jac}
\end{equation}
Thus $F_p$ is a Keller map and $A,B,C$ are algebraically independent.

\section{Inverse equation and exact recovery}

Equation \eqref{eq:phi} is equivalent to
\[
 \Phi(w)=wp(w)-cq(w).
\]
Eliminating $\gamma$ from \eqref{eq:PQ} yields
\begin{equation}
 H(T):=\Phi(T)-PT+cQ=0.                             \label{eq:inverse}
\end{equation}
Conversely, a generic root $w$ of \eqref{eq:inverse} recovers the source:
\begin{equation}
\begin{aligned}
 \gamma&=\frac{P-p(w)}c,&
 x&=\frac C\gamma,&
 u&=\frac w\gamma,&
 v&=u-1,\\
 y&=\frac vx,&
 t&=\frac{\gamma-1-av}{b},&
 z&=\frac t{x^2}.&
\end{aligned}                                      \label{eq:recovery}
\end{equation}
There is a particularly short certificate against an extraneous branch:
\begin{equation}
 H'(w)=p(w)-P=-c\gamma.                             \label{eq:derivative}
\end{equation}
A root with $\gamma=0$ would therefore be a repeated generic root.
As proved immediately below, $H$ is irreducible and hence separable in
characteristic zero.  Moreover, $C\ne0$ on a generic open set, so all
denominators in \eqref{eq:recovery} are valid at the generic point.

The target change is birational:
\[
 \C(A,B,C)=\C(P,Q,C),
\]
since $P=BC$, $Q=AC^2$, $B=P/C$, and $A=Q/C^2$.  If
$K=\C(A,B,C)$ and $L=\C(x,y,z)$, exact recovery gives
\begin{equation}
 L=K(w).                                            \label{eq:rootfield}
\end{equation}

If $\deg p=d-1$, then $\deg\Phi=d$.  The polynomial
$\Phi(T)-PT+cQ$ is irreducible over $\C(P,Q)$: it is irreducible in
$\C[P,Q,T]$ because it is primitive and linear in $Q$ with unit
coefficient, and it is primitive as a polynomial in $T$.  Gauss's lemma
applies.  Adjoining the independent transcendental $C$ preserves
irreducibility.  Hence \eqref{eq:rootfield} has degree $d$.

\section{Full symmetric monodromy}

We use the following classical theorem.

\begin{theorem}[Brink {\cite[Theorem 13]{Brink}}]\label{thm:brink}
Let $k$ be a field whose characteristic does not divide $d(d-1)$.  For
fixed $a_{d-1},\ldots,a_2\in k$, the polynomial
\[
 X^d+a_{d-1}X^{d-1}+\cdots+a_2X^2+SX+R
\]
has Galois group $S_d$ over $k(S,R)$.
\end{theorem}

Divide \eqref{eq:inverse} by the nonzero leading coefficient of $\Phi$.
Because $p(0)=0$, the fixed polynomial $\Phi$ has no linear term.  The
linear and constant coefficients of the normalized equation are independent
nonzero scalar multiples of $P$ and $Q$.  Taking $k=\C$ in
Theorem~\ref{thm:brink} gives
\begin{equation}
 \Gal\bigl(H/\C(P,Q)\bigr)=S_d.                     \label{eq:galois}
\end{equation}
The field $\C$ is algebraically closed, so this is already the geometric
group.  Adjoining the independent transcendental $C$ does not change the
splitting-field group.  Finally, \eqref{eq:rootfield} identifies its natural
root action with the action on the $d$ points of the generic fibre.
Therefore $\Mon(F_p)=S_d$, proving Theorem~\ref{thm:main}.

Gallagher's construction supplies admissible seeds in every degree
$d-1\ge2$.  Since $d>1$, the associated map cannot be a polynomial
automorphism.  This proves Corollary~\ref{cor:realization}.

\section{A compact rational family}

One particularly simple choice, useful for direct checking, is
\[
 b=c=1,\qquad
 p_d(w)=\frac{2w-dw^{d-1}}{d-2}.
\]
It satisfies \eqref{eq:seed}, has $\kappa=-(d+1)$, and gives
\[
 a=-\frac d{d-1},\qquad
 q_d(w)=\frac{w^2-(d-1)w^d}{d-2}.
\]
Thus
\[
 u=1+xy,\qquad
 \gamma=1-\frac d{d-1}xy+x^2z
\]
and
\begin{align*}
 A_d&=\frac{(d-2)u+u^2-(d-1)u^d\gamma^{d-2}}
              {(d-2)x^2},\\
 B_d&=\frac{(d-2)+2u-du^{d-1}\gamma^{d-2}}
              {(d-2)x},&
 C_d&=x\gamma.
\end{align*}
After $U=(d-2)P$ and $V=-(d-2)Q$, the inverse polynomial is simply
\[
 T^d-T^2+UT+V.
\]
This is also a direct instance of Brink's theorem.  Alternatively,
$T^d-T^2+UT$ is Morse over $\overline{\C(U)}$: its critical points are
simple, and equality of two critical values with ratio $r$ would imply
\[
 (d-2)r^d-dr^{d-1}+dr-(d-2)=0,
\]
forcing the transcendental parameter $U$ to be algebraic.  The classical
Morse-polynomial theorem \cite[Theorem 4.4.5]{Serre} again gives $S_d$.

There is an explicit collision.  Put
\[
 s_d=\frac{4-2^d}{d-2}.
\]
At the target $(A_d,B_d,C_d)=(s_d,s_d,1)$, the inverse polynomial has the
distinct roots $1$ and $2$.  Their recovery denominators are
\[
 \gamma(1)=\frac{d+2-2^d}{d-2}\ne0,\qquad
 \gamma(2)=2^{d-1}\ne0.
\]
Hence every displayed map is visibly noninjective.

\section{Priority and verification}

Gallagher \cite{Gallagher} proved the weighted-lift construction and the
occurrence of every generic degree at least three.  Brink
\cite{Brink} proved the Galois theorem used here.  Before this record,
MathOverflow answer 513470 and its linked Note~19
\cite{MO,Note19} proved $S_d$ for Gallagher's canonical tower for
$3\le d\le13$ and reported further finite checks.  We claim none of those
ingredients.  A source-specific search found no earlier statement that
Brink's theorem makes \emph{every} admissible Gallagher seed symmetric in
every degree.  Absence of a search hit is not proof of worldwide priority.

The accompanying exact verification has three parts.  A symbolic script
keeps all unconstrained seed coefficients symbolic and checks the endpoint,
integral, divisibility, elimination, and recovery identities for
$3\le d\le10$.  A separate PARI/GP resultant calculation verifies at the
specialization $U=1$ that all $d-1$ finite branch values are simple for
$3\le d\le20$, independently checking the geometric Morse mechanism.
Exact number-field specializations give arithmetic $S_d$ for
$3\le d\le10$; these are only arithmetic consistency checks and are not
presented as independent proof of geometric monodromy.

\medskip
\noindent\textbf{AI-assistance disclosure.}
This research record, proof organization, source search, and verification
code were produced with substantial assistance from OpenAI Codex agents.
The primary human researcher directed the program.  No person was contacted
by the agents.  The theorem passed an internal hostile-agent audit, which is
not peer review.

\begin{thebibliography}{9}

\bibitem{Brink}
D. Brink,
\emph{On alternating and symmetric groups as Galois groups},
Israel J. Math. \textbf{142} (2004), 47--60.
\href{https://doi.org/10.1007/BF02771527}{doi:10.1007/BF02771527}.

\bibitem{Gallagher}
A. Gallagher,
\emph{An infinite family of counterexamples to the Jacobian Conjecture in
dimension three: every generic fiber degree $n\ge3$ occurs},
Zenodo (20 July 2026).
\href{https://doi.org/10.5281/zenodo.21479195}{doi:10.5281/zenodo.21479195}.

\bibitem{MO}
\emph{Geometric degrees of counterexamples to the Jacobian conjecture in
dimension three}, MathOverflow answer 513470 (July 2026),
\url{https://mathoverflow.net/a/513470}.

\bibitem{Note19}
dasjoms,
\emph{Note 19---The Pin's Transposition}, public research note
(21 July 2026),
\url{https://github.com/dasjoms/jacobian-conjecture-counterexample-exploration/blob/main/jacobian_pin_transposition.md}.

\bibitem{Serre}
J.-P. Serre,
\emph{Topics in Galois Theory}, 2nd ed.,
Research Notes in Mathematics 1, A K Peters, 2008.

\end{thebibliography}

\end{document}
